\begin{abstract}
Structured state offers a way to bound a coding agent's active context
without carrying its full interaction history. Whether this also reduces
the cost of solving tasks depends on what the replacement memory preserves.
We evaluate SKILL.state-inspired context policies on Terminal-Bench,
varying state ownership, retained evidence, recent history, update protocol,
and schema. Model-written state reduces the median peak-prompt ratio to
\ModelStatePeakRatio{} on shared successes, yet raises the cost ratio to
\ModelStateCostRatio{} and lowers overall task success from
\ModelControlRate\% to \ModelStateRate\%. An exploratory comparison of
hybrid memory, which separates runtime records from model-written
interpretation, associates richer execution records with better success
and lower overhead than file references alone. Further enrichment does
not remove the full-history deficit. Retaining recent interactions gives
the clearest recovery: model-written and hybrid policies reach
\ModelHistorySixRate\% and \HybridHistorySixRate\%, against their respective
controls of \ModelControlRate\% and \HybridControlRate\%. The economic
outcomes differ: cost per observed success falls below the control for
model-written memory but remains above it for the hybrid variant.
The tested schema and protocol changes do not produce comparable recovery.
These results support evaluating structured memory through the evidence
it retains and the work needed to finish a task, rather than compression
alone. Recent-history retention is a measured partial remedy; the
incremental benefit of state over the same history window remains open.
\end{abstract}
